%! Change in Arithmetic and Geometric Sequences — Unit 4, Lesson 4.1 — Homework
\documentclass[10pt]{article}
\usepackage{precalculus-article}
\usepackage{precalculus-boxes}
\newcommand{\TallMath}[1]{$\displaystyle #1\rule[-1.4em]{0pt}{3.2em}$}

\begin{document}

\pageheader{Unit 4, Lesson 4.1}{Homework}
\namedateperiod

\vspace{0.10in}

\begin{enumerate}[itemsep=16pt]

\item The table below shows values of a sequence. Identify the sequence as arithmetic or
      geometric. Find $d$ or $r$, and write the explicit formula.

      \vspace{0.06in}
      \renewcommand{\arraystretch}{1.5}
      \begin{tabular}{|c|c|c|c|c|c|}
        \hline
        $n$    & 0  & 1  & 2   & 3   & 4   \\
        \hline
        $a(n)$ & 4  & 12 & 36  & 108 & 324 \\
        \hline
      \end{tabular}

      \vspace{0.06in}
      Type: \blank{3cm} \quad $d$ or $r$: \blank{2.5cm}

      Explicit formula: $a_n = $ \blank{5cm}

\item Two terms of a sequence are given: $a_2 = 7$ and $a_6 = 19$.

      \begin{enumerate}[label=(\alph*), itemsep=8pt]
        \item Assuming the sequence is arithmetic, find $d$.

              $d = \dfrac{a_6 - a_2}{6 - 2} = $ \blank{4cm}

        \item Write the explicit formula $a_n$.

              $a_n = $ \blank{6cm}

        \item Find $a_0$ and $a_{10}$.

              $a_0 = $ \blank{2.5cm} \quad $a_{10} = $ \blank{2.5cm}
      \end{enumerate}

\item A cell phone plan charges \$20 for month 0 (setup fee) and \$15 per month thereafter.

      \begin{enumerate}[label=(\alph*), itemsep=8pt]
        \item Write the explicit formula for the total cost $C_n$ after $n$ months.

              $C_n = $ \blank{5cm}

        \item Find the total cost after 6 months: $C_6 = $ \blank{3cm}

        \item In which month does the total cost first exceed \$200? \blank{3cm}
      \end{enumerate}

\item For the arithmetic sequence $a_n = 3 + 5n$, find the smallest value of $n$ such that
      $a_n > 250$.

      \vspace{0.06in}
      Setup: \blank{6cm}

      Smallest $n$: \blank{2.5cm} \quad Term value: $a_n = $ \blank{2.5cm}

\item \textbf{(AP-Style Multiple Choice)} A sequence is defined by $a_n = 4 + 3n$.
      Which of the following correctly describes this sequence?

      \medskip
      \begin{enumerate}[label=(\Alph*), itemsep=4pt]
        \item A geometric sequence with common ratio 3
        \item An arithmetic sequence with common difference 4
        \item An arithmetic sequence with common difference 3 and initial value 4
        \item A geometric sequence with initial value 4 and common ratio $\frac{1}{3}$
      \end{enumerate}

      \vspace{0.04in}
      Answer: \blank{1.5cm} \quad Explain:

      \vspace{0.04in}
      \writelines{2}

\item \textbf{(Extension)} Consider the sequence $1,\ 1,\ 2,\ 3,\ 5,\ 8,\ 13,\ \ldots$

      \begin{enumerate}[label=(\alph*), itemsep=8pt]
        \item Is this sequence arithmetic? Explain.

              \writeline

        \item Is this sequence geometric? Explain.

              \writeline

        \item What rule does this sequence follow? (Hint: look at pairs of consecutive terms.)

              \writelines{2}
      \end{enumerate}

\end{enumerate}

\end{document}
